\section{Where a specialist vendor can matter}

An aftermarket company cannot assume that every incomplete AI task is a market. A useful entry point has four properties. The task is valuable enough that completion, rather than a draft, is worth paying for. The customer can authorize a legitimate route. Success can be checked against an external state. And the work repeats often enough for the harness's setup and maintenance to be amortized. If any property is absent, the apparent opportunity may be a research curiosity, a consulting project, or a task that should remain manual.

The strongest initial domain may be narrow rather than glamorous. A vendor that knows one profession's forms, exceptions, account roles and acceptance checks can build a credible completion contract. A general consumer product must serve many domains with broadly safe defaults; the specialist can encode local vocabulary, customer workflows, and allowable tools. This is not evidence that the specialist will win. It is a reason to compare domain-specific harnesses against the current stock product and the customer's existing software, using a real task distribution rather than hand-picked demos.

Consider credential connection as one \emph{illustration}. An owner wants a recurring permitted workflow to reach several services. The model may know what data to move but the consumer product may lack a durable connection or may ask the user to re-enter context on every run. A specialist could coordinate official OAuth grants, approved service tokens, or an owner-operated browser handoff, then expose only task-scoped actions to the model. This is not a claim that a vault creates permission or that every website credential should be stored. The value hypothesis is narrower: once connection and revocation are handled, fewer user-active steps may be needed per accepted task. That hypothesis must be tested against the cost and risk of custody.

Account enrollment is a second illustration. A stock assistant may draft answers to a registration form but stop before the account becomes usable. A specialist harness might preserve form state, assemble owner-provided information, route identity checks and binding acceptance to the person, and verify the resulting account state. The person’s actions remain part of the task. The product claim would be that the owner spends less time reconstructing progress or repeating instructions, not that identity or platform controls disappear. Enrollment is worth pursuing only where the platform permits the route and the organization can accept the handoff.

Authorized defensive security work offers a third illustration without making security the paper's center. A defender may ask an assistant to summarize an alert, correlate approved telemetry, propose a bounded remediation, and verify whether the alert closed. A consumer product may stop at advice because it lacks the organization’s tools and approvals. A specialist product could connect read-only evidence, route a proposed action to the right operator, and record the result. This is a workflow hypothesis about legitimate defense, not a general authorization to perform security operations. It also illustrates why the unit sold is an accepted outcome rather than an eloquent recommendation.

These examples share a structure but not a single implementation. Each has an owner, an external system, missing execution context, a human or software checkpoint, and a verifiable terminal state. Other domains may be more attractive. The research task is to discover which part of the structure is repeatable across customers and which part demands expensive customization. A product that requires a bespoke integration for every customer can still be valuable, but its economics may resemble services more than software. The paper's superlative thesis depends on finding a repeatable wedge.

The specialist advantage is also constrained by distribution and trust. Frontier providers control the default interface and may add the missing tool or handoff. Existing vertical software vendors already have access to customer data and workflow budgets. A new harness vendor must show why its completion rate, user effort, evidence quality, or time to deploy is better. It must also be a good partner to the model provider and third-party platforms, because product behavior and permission rules can change. A moat based only on a prompt recipe is fragile; a moat based on a difficult but permitted acceptance contract may last longer, yet costs more to build.

\subsection*{The buyer is purchasing a terminal state}

The distinctive commercial offer is not that a customer can chat with a powerful model. They may already do that. The offer is that a recurring request can be submitted once, tracked without personal babysitting, and returned in one of a small number of legible states. Those states should be chosen with the buyer: accepted result, waiting on a specific external party, approval needed from a named role, or stopped for a documented reason. If the workflow cannot tell the difference between success and a plausible draft, it cannot sell dependable completion.

This framing changes which tasks deserve a first product. A spectacular one-off demonstration may attract attention but leave no repeatable demand. A mundane task with frequent recurrence, difficult context assembly, and a cheap independent acceptance check can be a better wedge. The relevant buyer may be an operations leader whose staff currently conducts the nudging, not the person who enjoys trying the newest model. Interviews should ask for recent task histories and failure traces, not only whether respondents are interested in an AI agent. Stated interest alone cannot establish willingness to pay or a tractable workflow.

One useful segmentation is by \emph{source of friction}. Some tasks are blocked by scattered context: the model has to see the right contract, policy, account state, or prior message. Others are blocked by a long sequence that a chat session will not hold together. Still others reach a nearly complete artifact but lack an external verifier. These segments suggest different minimum products. A retrieval layer may resolve the first; a durable controller may resolve the second; a domain acceptance test may resolve the third. Building all three before observing the dominant stop risks turning research into an expensive platform project.

The same segmentation also reveals when a vendor should say no. If the buyer cannot authorize a route, if the service disallows the action, or if the task's value depends on unrecorded human discretion, an aftermarket harness cannot promise programmatic completion. If the customer is already served by a stable integration, the correct competitor is that integration. A market map that treats every refusal or partial answer as addressable would inflate the opportunity and produce a fragile company. A narrower map, grounded in permission and acceptance, is more useful to founders deciding where to spend a year of engineering effort.

We therefore propose screening candidate markets with a small pilot protocol: define a representative task cohort; document the current user journey; record stock-product stop types and manual nudges; establish a deterministic baseline; and test whether a prototype harness raises verified completion at tolerable total cost. This pilot is a study design, not a promised return. Its most useful output may be a rejection of the market, especially when ordinary software already finishes the task or the supposed obstacle is a valid prohibition.
